% © 2026 Ing. David Jaroš — CC BY-NC-ND 4.0
%
% This work is licensed under a Creative Commons Attribution-NonCommercial-NoDerivatives
% 4.0 International License (CC BY-NC-ND 4.0).
%
% License History: Earlier drafts (up to v0.3) were released under CC BY 4.0.
% From v0.4 onward, all material is released under CC BY-NC-ND 4.0 to protect
% the integrity of the theoretical work during ongoing academic development.
%
% See LICENSE.md for full license text.

\ifdefined\INCLUDEMODE
\else
  \documentclass[12pt,a4paper]{article}
  \usepackage[a4paper, margin=2.5cm]{geometry}
  \usepackage{amsmath, amssymb, amsthm}
  \usepackage{mathtools}
  \usepackage{hyperref}
  \usepackage{xcolor}
  \usepackage{mdframed}
  \usepackage{booktabs}

  \newtheorem{theorem}{Theorem}[section]
  \newtheorem{lemma}[theorem]{Lemma}
  \newtheorem{proposition}[theorem]{Proposition}
  \newtheorem{conjecture}[theorem]{Conjecture}
  \theoremstyle{definition}
  \newtheorem{definition}[theorem]{Definition}
  \theoremstyle{remark}
  \newtheorem{remark}[theorem]{Remark}

  \newmdenv[backgroundcolor=yellow!20, linecolor=orange, linewidth=1.5pt]{gapbox}
  \newmdenv[backgroundcolor=green!15, linecolor=green!60!black, linewidth=1.5pt]{resultbox}
  \newmdenv[backgroundcolor=blue!8, linecolor=blue!50, linewidth=1.5pt]{insightbox}

  \title{\textbf{Electron Mass --- Step 1: Hecke Conditional Framework}\\
         \large If the Hecke Conjecture Holds, All Lepton Masses\\
         Reduce to One Free Parameter}
  \author{David Jaro\v{s}}
  \date{2026-03-06}

  \begin{document}
  \maketitle

  \begin{abstract}
  We analyse the structure of the lepton mass problem in UBT conditional
  on the Hecke eigenvalue conjecture.  If Hecke holds, the muon and tau
  masses are determined by Hecke eigenvalue \emph{ratios}, reducing the
  three independent Yukawa couplings of the Standard Model to a single
  free parameter: the electron mass scale $m_e$.
  This is already a major improvement over the SM (3 free $\to$ 1 free).
  The remaining question --- whether $m_e$ itself can be derived from
  first principles --- is the subject of Steps~2 and~3 of this series.
  \end{abstract}
\fi

%──────────────────────────────────────────────────────────────────────────────
\section{The Lepton Mass Problem in UBT}
\label{sec:setup}
%──────────────────────────────────────────────────────────────────────────────

The Standard Model lepton masses are:
\begin{align}
  m_e &= 0.510999\;\mathrm{MeV}, \\
  m_\mu &= 105.658\;\mathrm{MeV}, \\
  m_\tau &= 1776.86\;\mathrm{MeV}.
\end{align}
The ratios:
\begin{equation}
  \frac{m_\mu}{m_e} = 206.768, \qquad
  \frac{m_\tau}{m_e} = 3477.23.
\end{equation}

In the SM these require three independent Yukawa couplings $y_e, y_\mu, y_\tau$.
In UBT, the three generations are identified with $\psi$-winding modes
$n = 1, 2, 3$ (or the corresponding Hecke forms), but the mass
\emph{ratios} are not yet derived from first principles.

%──────────────────────────────────────────────────────────────────────────────
\section{The Hecke Conjecture}
\label{sec:hecke}
%──────────────────────────────────────────────────────────────────────────────

\begin{conjecture}[Hecke Eigenvalue Conjecture, Step~6]
\label{conj:hecke}
Let $f_1, f_2, f_3$ be the weight-$(2, 4, 6)$ Hecke newforms associated
with the three lepton generations, at a prime attractor level.  Then:
\begin{equation}
  \frac{m_\mu}{m_e} \;=\; \frac{a_{p^*}(f_2)}{a_{p^*}(f_1)}, \qquad
  \frac{m_\tau}{m_e} \;=\; \frac{a_{p^*}(f_3)}{a_{p^*}(f_1)},
\end{equation}
where $p^* \in \{137, 139\}$ is the prime attractor.
\end{conjecture}

\textbf{Numerical support} (from \texttt{step6\_hecke\_matches.tex}):
\begin{center}
\begin{tabular}{lccc}
  \toprule
  Prime & $m_\mu/m_e$ error & $m_\tau/m_e$ error & Status \\
  \midrule
  $p^* = 137$ & $0.37\%$ & $3.06\%$ & Numerical support \\
  $p^* = 139$ & $0.05\%$ & $1.63\%$ & Numerical support \\
  \bottomrule
\end{tabular}
\end{center}

%──────────────────────────────────────────────────────────────────────────────
\section{Reduction to One Free Parameter}
\label{sec:reduction}
%──────────────────────────────────────────────────────────────────────────────

\begin{proposition}[Conditional: 3 free $\to$ 1 free]
\label{prop:reduction}
Assume Conjecture~\ref{conj:hecke}.  Then:
\begin{equation}
  m_e, \; m_\mu, \; m_\tau
  \quad\text{are determined by}\quad
  m_e \;\text{alone.}
\end{equation}
Specifically:
\begin{equation}
  m_\mu = m_e \times r_{\mu/e}^{\mathrm{Hecke}}, \qquad
  m_\tau = m_e \times r_{\tau/e}^{\mathrm{Hecke}},
\end{equation}
where $r_{\mu/e}^{\mathrm{Hecke}}$ and $r_{\tau/e}^{\mathrm{Hecke}}$
are determined by Hecke eigenvalue ratios (no free parameters once the
LMFDB labels are fixed).
\end{proposition}

\begin{insightbox}
\textbf{Significance:}
The SM has 3 free Yukawa couplings for charged leptons.
Under Conjecture~\ref{conj:hecke}, UBT has 1 free parameter ($m_e$).
This is a reduction from 3 to 1 free parameter --- a major structural
improvement, even before $m_e$ itself is derived.
\end{insightbox}

%──────────────────────────────────────────────────────────────────────────────
\section{The Electron Mass: Gaps Y1 and Y2}
\label{sec:gaps}
%──────────────────────────────────────────────────────────────────────────────

The electron mass arises from the Higgs mechanism: $m_e = y_e \cdot v$
where $y_e$ is the Yukawa coupling and $v = 246\;\mathrm{GeV}$ is the
Higgs VEV.  In UBT:

\begin{description}
  \item[Gap Y1] Derive $y_e$ from $S[\Theta]$.
  \item[Gap Y2] Derive $v$ from the effective potential $V_{\mathrm{eff}}(\theta_0)$
                on the $\psi$-circle without $m_e$ as input.
\end{description}

Both gaps are currently open.  Steps~2 and~3 of this series explore two
independent approaches.

\begin{remark}[Dimensional analysis]
$m_e \sim \hbar/(c\,R_\psi)$ where $R_\psi$ is the $\psi$-circle radius.
The current calibration is $R_\psi = \hbar/(m_e c)$, which is circular.
A non-circular derivation requires fixing $R_\psi$ from a different
condition (see Step~2).
\end{remark}

%──────────────────────────────────────────────────────────────────────────────
\section{Conclusion}
\label{sec:conclusion}
%──────────────────────────────────────────────────────────────────────────────

\begin{resultbox}
\textbf{Conditional result (LABEL: conditional on Hecke):}
\begin{itemize}
  \item IF Conjecture~\ref{conj:hecke} is proved:
        all lepton masses are determined by $m_e$ alone.
        UBT has 1 free parameter in the lepton sector vs.\ 3 in SM.
  \item The Hecke conjecture currently has numerical support at $0.05\%$--$3.06\%$
        level but is not proved.
  \item The absolute scale $m_e$ requires independent derivation (Gaps Y1, Y2).
\end{itemize}

\textbf{Status of Y1, Y2:} Open.  See
\texttt{step2\_geometric\_R\_psi.tex} and
\texttt{step3\_modular\_period.tex} for approaches.
\end{resultbox}

\ifdefined\INCLUDEMODE
\else
  \end{document}
\fi
